% Part: lambda-calculus
% Chapter: introduction
% Section: composition

\documentclass[../../../include/open-logic-section]{subfiles}

\begin{document}

\olfileid{lam}{int}{com}
\olsection{توابع \usetoken{S}{lambda definable} تحت ترکیب بسته‌اند}

\begin{lem}
توابع !!{lambda definable} تحت ترکیب بسته‌اند.
\end{lem}

\begin{proof}
فرض کنید~$f$ با ترکیب از~$h$، $g_0$،
\dots،~$g_{k-1}$ تعریف شده باشد. با فرض اینکه~$h$، $g_0$،
\dots،~$g_{k-1}$ به‌ترتیب به‌وسیلهٔ~$H$، $G_0$،
\dots،~$G_{k-1}$ !!{lambda defined} باشند، باید ترم~$F$ای بیابیم که~$f$
را !!{lambda define}s. اما می‌توانیم~$F$ را به‌سادگی چنین تعریف کنیم:
\[
F(x_0, \dots, x_{l-1}) = H(G_0(x_0, \dots, x_{l-1}),
\dots, G_{k-1}(x_0, \dots, x_{l-1})).
\]
به بیان دیگر، زبان حساب لامبدا برای بازنمایی ترکیب کاملاً مناسب است.
\end{proof}

\end{document}
